\begin{code}[hide]%
\>[0]\AgdaKeyword{open}\AgdaSpace{}%
\AgdaKeyword{import}\AgdaSpace{}%
\AgdaModule{Data.Product}\<%
\\
\>[0]\AgdaKeyword{open}\AgdaSpace{}%
\AgdaKeyword{import}\AgdaSpace{}%
\AgdaModule{Data.Fin}\<%
\\
\>[0]\AgdaKeyword{open}\AgdaSpace{}%
\AgdaKeyword{import}\AgdaSpace{}%
\AgdaModule{Data.Nat}\<%
\end{code}

\section{\label{sec:related}Related work}

% Containers
The array theory presented herein draws inspiration from earlier work on
containers~\citep{containers}.  We restrict the shapes of containers to
products of natural numbers, and our reshapes correspond to Cartesian
morphisms~\citep{10.1007/3-540-36576-1_2}.  Containers prove extremely useful
when modelling arrays.  On the one hand, an explicit notion of shapes and
indices affords a straightforward representation of arrays in concrete memory ---
we merely need to translate indices of a given shape into offsets.  On the
other hand, because arrays are modelled as functions, reasoning about array
expressions becomes simpler; we can frequently exploit properties such as
map-fusion through normalisation.  Further developments of arrays as containers
can be found in \citet{sinkarovs_2020}.  As this work deals with finite arrays,
the theory of combinatorial species~\citep{Yorgey2014}---which are essentially
finitary containers---is also relevant.  Moreover, since our shapes possess a
monoidal structure, the theory of operads~\citep{May1972} can be employed to
interpret the proposed arrays.

% Examples of other examples that can be modeled in a similar style.
Due to limited space we presented only a single running example that
demonstrates the application of our technique.  Nevertheless, a comparable
treatment of arrays has proved valuable in the following two scenarios.
In \citet{10.1145/3609024.3609410}, we consider an algorithm for blocked matrix
multiplication.  Much as with scan, the tree structure of array shapes enables
the encoding of a blocking strategy for the input matrices.  We define a
high-level algorithm for blocked matrix multiplication, prove its correctness,
and use it as a guide to derive an implementation in the array language SaC
that performs on par with state-of-the-art libraries such as OpenBLAS and MKL.
In \citet{10.1145/3747524}, we employ a similar array framework to define
convolutional neural networks using high-level array combinators.  We provide
an automatic-differentiation facility in Agda together with a
semantics-preserving optimiser and a code generator targeting Futhark for the
language of array combinators.

% Conal's work
Many ideas that appear in \citet{Elliott:scan} resonate with the spirit of this
paper.  \citeauthor{Elliott:scan} defines a language of (generic) functors
comprising zero, one, and singleton, closed under sums, products and
composition.  For a given algorithm such as scan, one then defines an instance
of the algorithm for every case of that functor language.  Consequently, the
chosen algorithm can be computed for data structures expressible as generic
functors.  Containers are isomorphic~\citep{ABBOTT20053} to strictly positive
functors, so all generic operations can be defined on containers as well.  The
key differences in the present work are threefold:  
%
(i)  The dependently-typed framework permits reasoning about rank polymorphism
--- knowledge of array shapes may be deferred until run-time --- whereas the
decomposition into generic functors advocated by \citeauthor{Elliott:scan} must
occur statically;  
%
(ii)  Here we abstract over array representation --- by selecting a model and
combinators the algorithm is guaranteed to compute on the chosen representation
--- whereas computation via generic functors proceeds by translation, leaving
optimisation to the compiler and thereby introducing separate practical
challenges;  
%
(iii)  Most importantly, we furnish a formal proof that isomorphic
representations (as given by reshapes) yield identical results, whereas with
generic functors one must manually ensure, for example, that scans over
\texttt{Pair × Pair} and \texttt{Pair ∘ Pair} coincide.
Finally, our use of monads provides precise numbers that can be taken
into account when adapting specifications to a particular hardware
model.

% Array languages
There exists a substantial body of array languages that support rank
polymorphism to varying extents.  Early untyped languages such as
APL~\citep{Iverson:1962:APL}, J~\citep{jlang}, and Nial~\citep{nial} provide
run-time rank polymorphism without any type or conformance guarantees.  They
offer a rich repertoire of rank-polymorphic array combinators and are valuable
for rapid prototyping.  Within this tradition, Trenchard More initiated the
development of array theories~\citep{More73} that focus on identifying
universal equalities in the context of built-in rank-polymorphic combinators.
This line of work culminated in Lenore Mullin's
$\Psi$-calculus~\citep{Mullin88}, which is based on a universal abstraction of
iteration spaces akin to the \AF{imap} operator from Section~\ref{sec:arrop}.

Inspired by these ideas, many functional array languages strive to support rank
polymorphism: SaC~\citep{sac}, Qube~\citep{Qube}, Accelerate~\citep{MCGN15},
Jeremy Gibbons' Naperian Functors~\citep{Gibbons:APLicative},
Remora~\citep{Remora}, and ATL~\citep{ATL}.  While languages such as SaC or
Accelerate primarily target the generation of efficiently executable code, the
others are chiefly concerned with providing static guarantees against
out-of-bounds errors.  A fully generic setting in which ranks may only be known
at run-time necessitates dependent types, as embodied by the array theory
presented in this paper or by the type system of Qube.  When ranks are
restricted to static (compile-time) values, many static guarantees can be
obtained with a less expressive type system.  For instance,
\citet{Gibbons:APLicative} offers a Haskell framework --- similar in spirit to
our work --- for handling statically-ranked multidimensional arrays.
Remora~\citep{Remora} proposes a novel type system (without full dependent
types) for statically-ranked arrays, furnishing numerous APL-specific features
such as shape lifting.  ATL~\citep{ATL} is a statically-ranked Coq-embedded
language that provides safe static indexing.  In principle, all of the
aforementioned languages could implement a static variant of the scan
specification presented here.  SaC, for example,
demonstrates~\citep{10.1145/3520306.3534500} the expressive power of rank
polymorphism.  Nonetheless, reproducing our equivalence proof based on reshapes
within those systems would likely be challenging.

Futhark~\citep{henriksen:phdthesis} is a functional array language that
prioritises parallel performance and type safety.  In contrast to the languages
discussed above, it does not support rank polymorphism at all, but it does
provide a built-in scan primitive.  Although convenient for users, this
primitive often renders it difficult or impossible to control the run-time
behaviour of a particular application.


%\item Hinze Scan Algebra
\citet{Hinze:ScanAlgebra} explores the compositional nature of
the in-place algorithm introduced by \citet{GEBlelloch:prefixsums}.  It would be interesting to
investigate whether a tidy translation exists between our specification and the
one presented in that paper.

%\item Extraction
\citet{abs-2105-10819} demonstrate that specifications from a framework akin
to the one presented here can be translated into high-performance languages.
This technique, termed ``extraction'', leverages the reflection capabilities of
the underlying dependently-typed language.  With a modest amount of effort,
Agda specifications can be compiled into an array language of the user's
choice.

% \item CwFs
Our encoding of arrays bears a strong resemblance to the encoding of type
theories using Categories with Families (CwFs)~\citep{Castellan2021}.  This is
unsurprising, since arrays are essentially functions: \AF{imap} behaves like
$\lambda$, and \AF{\_[\_]} (selection) behaves like function application.
Accordingly, the analogy proceeds as follows:
\begin{code}[hide]%
\>[0]\AgdaKeyword{postulate}\<%
\\
\>[0][@{}l@{\AgdaIndent{0}}]%
\>[2]\AgdaPostulate{Ty}%
\>[7]\AgdaSymbol{:}\AgdaSpace{}%
\AgdaPrimitive{Set}\<%
\\
\>[0]\AgdaKeyword{postulate}\<%
\end{code}
\begin{mathpar}
\codeblock{\begin{code}%
\>[0][@{}l@{\AgdaIndent{1}}]%
\>[2]\AgdaPostulate{Con}%
\>[7]\AgdaSymbol{:}\AgdaSpace{}%
\AgdaPrimitive{Set}\<%
\\
%
\>[2]\AgdaPostulate{Sub}%
\>[7]\AgdaSymbol{:}\AgdaSpace{}%
\AgdaPostulate{Con}\AgdaSpace{}%
\AgdaSymbol{→}\AgdaSpace{}%
\AgdaPostulate{Con}\AgdaSpace{}%
\AgdaSymbol{→}%
\>[22]\AgdaPrimitive{Set}\<%
\\
%
\>[2]\AgdaPostulate{Tm}%
\>[7]\AgdaSymbol{:}\AgdaSpace{}%
\AgdaPostulate{Con}\AgdaSpace{}%
\AgdaSymbol{→}\AgdaSpace{}%
\AgdaPostulate{Ty}\AgdaSpace{}%
\AgdaSymbol{→}\AgdaSpace{}%
\AgdaPrimitive{Set}\<%
\end{code}}
\and
\codeblock{\begin{code}%
%
\>[2]\AgdaPostulate{Sh}%
\>[11]\AgdaSymbol{:}\AgdaSpace{}%
\AgdaPrimitive{Set}\<%
\\
%
\>[2]\AgdaPostulate{Reshape}%
\>[11]\AgdaSymbol{:}\AgdaSpace{}%
\AgdaPostulate{Sh}\AgdaSpace{}%
\AgdaSymbol{→}\AgdaSpace{}%
\AgdaPostulate{Sh}\AgdaSpace{}%
\AgdaSymbol{→}\AgdaSpace{}%
\AgdaPrimitive{Set}\<%
\\
%
\>[2]\AgdaPostulate{Ar}%
\>[11]\AgdaSymbol{:}\AgdaSpace{}%
\AgdaPostulate{Sh}\AgdaSpace{}%
\AgdaSymbol{→}\AgdaSpace{}%
\AgdaPrimitive{Set}\AgdaSpace{}%
\AgdaSymbol{→}\AgdaSpace{}%
\AgdaPrimitive{Set}\<%
\end{code}}
\end{mathpar}

Array shapes play the role of contexts, and array indices correspond to
variables.  Reshapes act as substitutions that permit the remapping of
variables/indices within an array expression.  Arrays correspond to terms,
with one difference that is visible in the signatures above: terms inhabit a
closed universe of object-level types \AD{Ty}, whereas array elements may be
of any Agda type (\AD{Set}).  The crucial insight for both systems is that substitutions
possess an inductive structure, allowing the action of a substitution (\ie{}
reshaping) to be defined inductively as well.

In a sense, the proposed array theory constitutes a rudimentary type theory,
wherein the only types that variables may assume in a context (shape) are of
the form \AD{Fin}\,\AB{n} for some natural number \AB{n}.  It would be
intriguing to develop a closed array theory in which indices and arrays inhabit
a single universe.  In such a setting, array shapes would themselves be arrays,
guaranteeing that any operation on shapes is immediately applicable to arrays.
A minimal prerequisite for this construction would be the ability to embed the
array theory presented in this paper into such a type theory.
Our experiments, however, suggest that
attaining this goal might require something akin to very dependent
types~\citep{altenkirch_et_al:LIPIcs.TYPES.2022.10}.
